\documentclass[tcc,capa]{texifsul}
\setlength{\parindent}{2,96382876979577em}

\usepackage[utf8]{inputenc} % acentuacao
\usepackage{graphicx} % para inserir figuras
\usepackage[T1]{fontenc}
\usepackage{multirow}
\usepackage{listings}
\usepackage{xcolor}

\hypersetup{
    hidelinks,
    unicode=true,
    linktoc=all
}

\unidade{Câmpus Gravataí}

\nomecurso{Superior de Tecnologia em Análise e Desenvolvimento de Sistemas}
\titulocurso{Tecnólogo em Análise e Desenvolvimento de Sistemas}

\unidadeeng{Campus Gravataí}

\title{CRAWLER AI: COMPARAÇÃO ENTRE EXTRAÇÃO DE DADOS VIA MODELOS DE LINGUAGEM E MÉTODOS TRADICIONAIS DE WEB SCRAPING}

\author{}{Felipe Machado Schultz}

\advisor[Prof.~Titulação.]{}{Fernando Abraão Afonso}

\banca[Prof.~Titulação.]{}{Nome membro da banca}
\banca[Prof.~Titulação.]{}{Nome membro da banca}

%Palavras-chave em PT_BR
\keyword{Web Scraping}
\keyword{Modelos de Linguagem}
\keyword{Extração de Dados}
\keyword{Automação Web}
\keyword{Inteligência Artificial}

%Palavras-chave em EN_US
\keywordeng{Web Scraping}
\keywordeng{Large Language Models}
\keywordeng{Data Extraction}
\keywordeng{Web Automation}
\keywordeng{Artificial Intelligence}

\begin{document}

\maketitle

\sloppy

\begin{aprovacao}{22 de junho de 2026}
\end{aprovacao}

% \begin{agradecimentos}
%   Agradeço\ldots
% \end{agradecimentos}

% \begin{epigrafe}
%   The question of whether a computer can think is no more interesting than the question of whether a submarine can swim.\\
%   {\sc --- Edsger W. Dijkstra}
% \end{epigrafe}

\begin{abstract}
\noindent
% TODO: Escrever o resumo após a conclusão do trabalho (~300 palavras).
% Deve conter: contexto, problema, objetivo, metodologia, resultados e conclusão.
\end{abstract}

\begin{englishabstract}
\noindent
% TODO: Escrever o abstract em inglês após a conclusão do trabalho (~300 palavras).
\end{englishabstract}

\listoffigures
\listoftables

\begin{listofabbrv}{RabbitMQ}
    \item[ADS] Análise e Desenvolvimento de Sistemas
    \item[API] Application Programming Interface
    \item[CSS] Cascading Style Sheets
    \item[DOM] Document Object Model
    \item[GPT] Generative Pre-trained Transformer
    \item[HTML] HyperText Markup Language
    \item[HTTP] HyperText Transfer Protocol
    \item[IA] Inteligência Artificial
    \item[IFSul] Instituto Federal de Educação, Ciência e Tecnologia Sul-rio-grandense
    \item[JSON] JavaScript Object Notation
    \item[LLM] Large Language Model
    \item[NLP] Natural Language Processing
    \item[ORM] Object-Relational Mapping
    \item[PLN] Processamento de Linguagem Natural
    \item[RabbitMQ] Rabbit Message Queue
    \item[REST] Representational State Transfer
    \item[SPA] Single Page Application
    \item[SQL] Structured Query Language
    \item[TCC] Trabalho de Conclusão de Curso
    \item[URL] Uniform Resource Locator
    \item[XPath] XML Path Language
\end{listofabbrv}

\tableofcontents

% ============================================================
\chapter{Introdução}
% ============================================================

A internet é hoje o maior repositório de informações já construído pela humanidade. Segundo a empresa de análise de dados Statista, o volume de dados gerados e consumidos globalmente ultrapassa dezenas de zettabytes por ano, com projeções de crescimento contínuo impulsionadas pela expansão de serviços digitais, redes sociais, plataformas de comércio eletrônico e sistemas governamentais abertos \cite{Statista2024}. Nesse contexto, a capacidade de coletar, estruturar e analisar informações disponíveis publicamente na web tornou-se uma habilidade estratégica para empresas, pesquisadores e desenvolvedores.

O web scraping — técnica de extração automatizada de dados de páginas web — surgiu como resposta a essa demanda. Por meio de programas que simulam o acesso humano a páginas, é possível coletar e estruturar grandes volumes de informação de forma automatizada, alimentando pipelines de dados, sistemas de monitoramento de preços, ferramentas de análise de mercado e diversas outras aplicações \cite{Mitchell2018}. Técnicas como a análise do HTML das páginas por meio de seletores CSS e expressões XPath, a utilização de bibliotecas como o BeautifulSoup e frameworks como o Scrapy consolidaram-se ao longo dos anos como as abordagens dominantes para esse tipo de tarefa.

No entanto, o scraping tradicional apresenta uma limitação estrutural fundamental: ele é intrinsecamente dependente da forma como o conteúdo está organizado no código-fonte da página. Qualquer alteração no layout ou na estrutura do HTML — uma prática comum em aplicações web modernas — invalida imediatamente os seletores e regras definidos, exigindo manutenção manual constante dos coletores. Essa fragilidade representa um custo operacional significativo, especialmente em cenários de monitoramento contínuo de múltiplas fontes de dados \cite{GlezPena2014}.

Paralelamente, a última década testemunhou uma transformação radical no campo da Inteligência Artificial com o advento dos Modelos de Linguagem de Grande Escala (LLMs, do inglês \textit{Large Language Models}). Modelos como o GPT-4 da OpenAI demonstraram capacidades notáveis de compreensão e geração de texto, incluindo a habilidade de interpretar documentos estruturados, extrair informações específicas e produzir saídas em formatos predefinidos como o JSON \cite{OpenAI2023}. Essa capacidade de entendimento semântico — não dependente da estrutura sintática do documento — abre uma nova perspectiva para o problema da extração de dados web.

Um sistema de scraping baseado em LLM, em teoria, seria capaz de extrair informações de uma página web com base em uma descrição em linguagem natural do que deve ser coletado, sem necessidade de conhecer previamente a estrutura do HTML. Em vez de regras frágeis baseadas em seletores, o modelo recebe o conteúdo textual da página e um esquema de dados desejado, retornando as informações estruturadas. Essa abordagem promete maior resiliência a mudanças de layout e menor custo de manutenção, mas introduz novos desafios relacionados a custo financeiro por token, latência de resposta e dependência de serviços externos.

Diante desse cenário, o presente trabalho propõe o desenvolvimento do sistema \textbf{Crawler AI}, uma aplicação de web scraping que utiliza LLMs para a extração de dados estruturados, e sua comparação sistemática com a abordagem tradicional baseada em seletores CSS/XPath e BeautifulSoup. O objetivo central é fornecer uma análise empírica das vantagens, limitações e casos de uso mais adequados para cada abordagem, contribuindo para o embasamento técnico de decisões de engenharia de software nesse domínio.

Este trabalho está organizado da seguinte forma: o Capítulo~\ref{cap:justificativa} apresenta a justificativa e motivação do estudo; o Capítulo~\ref{cap:objetivos} enuncia os objetivos geral e específicos; o Capítulo~\ref{cap:revisao} traz a revisão de literatura, cobrindo fundamentos de web scraping e modelos de linguagem; o Capítulo~\ref{cap:metodologia} descreve a arquitetura do sistema implementado e a metodologia de comparação adotada; o Capítulo~\ref{cap:resultados} apresenta e discute os resultados obtidos nos experimentos; e o Capítulo~\ref{cap:conclusao} reúne as conclusões e propostas de trabalhos futuros.

% ============================================================
\chapter{Justificativa}
\label{cap:justificativa}
% ============================================================

A extração automatizada de dados da web é uma necessidade crescente em contextos empresariais e acadêmicos. Monitoramento de preços em e-commerces, coleta de dados para treinamento de modelos de aprendizado de máquina, análise de sentimentos em redes sociais e agregação de notícias são apenas alguns exemplos de aplicações que dependem de scrapers confiáveis e eficientes. No contexto do curso de Análise e Desenvolvimento de Sistemas, compreender e avaliar as ferramentas e abordagens disponíveis para esse tipo de tarefa é uma competência diretamente aplicável ao mercado de trabalho.

As abordagens tradicionais de web scraping, embora consolidadas, enfrentam desafios práticos significativos. O principal deles é a fragilidade dos coletores frente a mudanças na estrutura das páginas. Um scraper baseado em seletores CSS que funciona perfeitamente hoje pode falhar completamente amanhã caso o time de desenvolvimento do site altere o nome de uma classe CSS ou reestruture o DOM. Esse comportamento demanda monitoramento constante e manutenção reativa, gerando custo operacional que muitas vezes não é contabilizado na fase de projeto.

Além disso, a crescente adoção de Aplicações de Página Única (SPAs, do inglês \textit{Single Page Applications}) — construídas com frameworks como React, Angular e Vue — adiciona uma camada extra de complexidade: o conteúdo da página é renderizado dinamicamente via JavaScript, tornando inacessível ao scraping via requisição HTTP simples. Para esses casos, é necessário recorrer à automação de navegadores reais (como o Playwright ou o Selenium), o que aumenta a complexidade de infraestrutura e o tempo de execução.

Os LLMs representam uma abordagem conceitualmente diferente. Em vez de descrever onde está a informação (seletor), o desenvolvedor descreve o que é a informação (esquema semântico). O modelo, treinado em quantidades massivas de texto da internet, possui capacidade de inferir o contexto e extrair os dados relevantes mesmo que a estrutura da página mude. Essa propriedade tem o potencial de reduzir drasticamente o custo de manutenção de scrapers, uma vez que a lógica de extração é expressa em termos semânticos em vez de estruturais.

Contudo, a adoção de LLMs para scraping não é isenta de custos e riscos. O uso de APIs de modelos como o GPT-4 da OpenAI implica custos financeiros por volume de tokens processados, além de latência de rede e possível indisponibilidade do serviço. Há também o risco de alucinações — casos em que o modelo infere dados que não estão presentes na página. Questões de privacidade e conformidade com termos de uso de serviços terceiros também precisam ser consideradas.

Apesar da relevância prática do tema, a literatura acadêmica carece de estudos comparativos formais e sistemáticos entre as duas abordagens. A maioria das comparações disponíveis é feita em blogs e artigos técnicos não revisados por pares, sem uma metodologia de avaliação rigorosa. Este trabalho busca preencher essa lacuna ao implementar ambas as abordagens sobre os mesmos conjuntos de dados e avaliá-las com métricas objetivas, produzindo evidências empíricas que possam orientar escolhas de engenharia de software em projetos reais.

% ============================================================
\chapter{Objetivos}
\label{cap:objetivos}
% ============================================================

\section{Objetivo Geral}

Implementar o sistema Crawler AI, uma aplicação de web scraping baseada em Modelos de Linguagem de Grande Escala, e compará-lo sistematicamente com a abordagem tradicional de extração de dados baseada em seletores CSS/XPath, avaliando as duas abordagens quanto à acurácia dos dados extraídos, ao tempo de processamento, ao custo operacional e ao esforço de manutenção.

\section{Objetivos Específicos}

\begin{enumerate}
    \item Implementar o pipeline de scraping via LLM do sistema Crawler AI, composto pelos estágios de busca (\textit{fetch}), limpeza (\textit{clean}) e extração (\textit{extract}) utilizando o modelo GPT-4o mini da OpenAI;
    \item Implementar um módulo de scraping tradicional como \textit{baseline} de comparação, utilizando a biblioteca BeautifulSoup com seletores CSS e expressões XPath;
    \item Definir um conjunto de métricas de avaliação para a comparação entre as abordagens, incluindo acurácia, tempo de resposta, custo por extração e esforço de implementação e manutenção;
    \item Selecionar e executar experimentos com cenários representativos de casos de uso reais, contemplando diferentes tipos de páginas web (e-commerce, portais de notícias e páginas com dados tabulares);
    \item Analisar e comparar os resultados obtidos em cada cenário, identificando os pontos fortes e limitações de cada abordagem;
    \item Documentar os resultados e conclusões de forma a contribuir com subsídios técnicos para a escolha entre abordagens de scraping em projetos de desenvolvimento de sistemas.
\end{enumerate}

% ============================================================
\chapter{Revisão de Literatura}
\label{cap:revisao}
% ============================================================

% TODO: Esta seção deve ter entre 3.000 e 4.500 palavras.
% Referenciar ao longo do texto com \cite{}.

\section{Web Scraping: Conceito e Evolução}

% TODO: Escrever ~400 palavras cobrindo:
% - Definição de web scraping
% - Histórico: primeiros crawlers (anos 90), evolução com AJAX, SPAs
% - Casos de uso: pesquisa acadêmica, monitoramento, BI, ML datasets
% - Aspectos legais e éticos (robots.txt, termos de uso)
% Referências sugeridas: Mitchell2018, GlezPena2014

\section{Abordagens Tradicionais de Web Scraping}

% TODO: Escrever ~600 palavras cobrindo:
% - Parsing de HTML com BeautifulSoup e lxml
% - Seletores CSS e expressões XPath: sintaxe e uso
% - Frameworks completos: Scrapy (arquitetura de spider)
% - Automação de navegador: Selenium, Playwright (motivação: SPAs e JS)
% - Vantagens: controle total, sem custo de API, determinístico
% - Desvantagens: fragilidade a mudanças de estrutura, alto custo de manutenção
% Referências sugeridas: Mitchell2018, GlezPena2014, Playwright2024, BeautifulSoup2024

\section{Modelos de Linguagem de Grande Escala}

% TODO: Escrever ~600 palavras cobrindo:
% - Definição de LLM
% - Arquitetura Transformer (Vaswani et al., 2017): mecanismo de atenção
% - Pré-treinamento e fine-tuning: GPT, BERT e variantes
% - GPT-3 (Brown et al., 2020) e GPT-4 (OpenAI, 2023): capacidades emergentes
% - Conceito de tokens, context window e custo por uso
% Referências sugeridas: Vaswani2017, Brown2020, OpenAI2023, Zhao2023

\section{LLMs para Extração de Informação Estruturada}

% TODO: Escrever ~500 palavras cobrindo:
% - Information Extraction (IE) como tarefa de NLP
% - LLMs como extratores zero-shot e few-shot
% - JSON mode / structured output: garantia de formato de resposta
% - Exemplos de aplicações: extração de entidades, preenchimento de formulários, parsing de documentos
% - Limitações: alucinação, custo, latência
% Referências sugeridas: Brown2020, OpenAI2023, Wei2022

\section{Engenharia de Prompts}

% TODO: Escrever ~400 palavras cobrindo:
% - Definição e importância do prompt engineering
% - Técnicas: zero-shot, few-shot, chain-of-thought (Wei et al., 2022)
% - Impacto na qualidade da extração
% - Boas práticas: instrução clara, exemplos de entrada/saída, restrições explícitas
% Referências sugeridas: Wei2022, Zhao2023

\section{Trabalhos Relacionados}

% TODO: Escrever ~500 palavras cobrindo:
% - Pesquisar e citar estudos existentes que comparam scraping tradicional e LLM
% - Ferramentas como ScrapeGraphAI, Firecrawl, Jina AI Reader
% - Artigos que usam LLMs para extração de dados de páginas web
% - Lacunas na literatura que este trabalho pretende preencher

% ============================================================
\chapter{Metodologia}
\label{cap:metodologia}
% ============================================================

\section{Visão Geral do Sistema Crawler AI}

O Crawler AI é um sistema web desenvolvido para executar tarefas de extração de dados de páginas da internet utilizando Modelos de Linguagem. O sistema permite que o usuário autenticado informe uma ou mais URLs, descreva o esquema de dados desejado em formato JSON e, opcionalmente, forneça um contexto adicional em linguagem natural. A extração é então realizada de forma assíncrona, com os resultados disponibilizados na interface assim que processados.

A arquitetura do sistema segue o padrão de processamento assíncrono baseado em filas de mensagens, separando a camada de apresentação (\textit{frontend}) do processamento pesado de extração (\textit{worker}). Essa separação garante que a interface permaneça responsiva mesmo durante o processamento de múltiplas URLs simultaneamente.

\section{Arquitetura do Sistema}

% TODO: Inserir figura com diagrama da arquitetura completa.
% Descrever os componentes:

O sistema é composto por três camadas principais que se comunicam por meio de filas de mensagens:

\subsection{Frontend (Interface Web)}

A interface do usuário foi desenvolvida com React 19 e o framework TanStack Start, utilizando renderização no servidor (SSR) via Nitro. O roteamento é gerenciado pelo TanStack Router, que adota uma abordagem de roteamento baseada em arquivos. Para o gerenciamento de estado assíncrono e cache de dados, utiliza-se o TanStack Query. A autenticação de usuários é implementada com a biblioteca Better Auth, com sessões persistidas em banco de dados PostgreSQL. O ambiente de execução e gerenciador de pacotes adotado é o Bun.

\subsection{Worker de Crawling (Python)}

O componente responsável pelo processamento é um worker desenvolvido em Python que consome mensagens de uma fila RabbitMQ. Ao receber uma mensagem contendo a URL e o esquema de extração, o worker executa o pipeline de três estágios descrito na Seção~\ref{sec:pipeline} e publica o resultado em outra fila para persistência.

\subsection{Infraestrutura de Suporte}

A comunicação entre o frontend e o worker é mediada pelo RabbitMQ, um broker de mensagens que implementa o protocolo AMQP. O banco de dados relacional PostgreSQL é utilizado para persistência de usuários, sessões, jobs de extração e seus respectivos resultados. Toda a infraestrutura é provisionada via Docker Compose.

\section{Pipeline de Extração via LLM}
\label{sec:pipeline}

O núcleo do sistema Crawler AI é um pipeline de processamento sequencial composto por três estágios: busca (\textit{fetch}), limpeza (\textit{clean}) e extração (\textit{extract}). Cada estágio é implementado como um módulo independente, permitindo substituição ou extensão de qualquer componente sem afetar os demais.

\subsection{Estágio 1: Busca (\textit{Fetch})}

O estágio de busca é responsável por obter o conteúdo HTML bruto da página alvo. O sistema adota uma estratégia de dois níveis para maximizar a taxa de sucesso:

\begin{enumerate}
    \item \textbf{Requisição HTTP direta}: inicialmente, o sistema realiza uma requisição HTTP utilizando a biblioteca \texttt{requests}, com cabeçalhos que simulam um navegador real (incluindo User-Agent do Chrome e preferência de idioma português). Esta abordagem é rápida e de baixo custo computacional;
    \item \textbf{Automação via Playwright}: caso a requisição HTTP falhe (por bloqueio, redirecionamento ou conteúdo renderizado via JavaScript) ou retorne conteúdo insuficiente — detectado por marcadores característicos de SPAs —, o sistema recorre ao Playwright com Chromium em modo headless. O navegador realiza o carregamento completo da página, incluindo a execução de JavaScript, aguardando dois segundos adicionais após o evento de \textit{load} para garantir que o conteúdo dinâmico seja renderizado.
\end{enumerate}

Os timeouts são configuráveis via variáveis de ambiente: 15 segundos para HTTP e 60 segundos para Playwright. O módulo trata explicitamente os códigos de erro HTTP 403 (acesso negado), 429 (muitas requisições) e 503 (serviço indisponível).

\subsection{Estágio 2: Limpeza (\textit{Clean})}

O HTML bruto obtido no estágio anterior contém uma quantidade significativa de ruído que não contribui para a extração semântica: tags de script, folhas de estilo, elementos de navegação, rodapés, iframes, elementos SVG, metadados e outros. O estágio de limpeza utiliza a biblioteca BeautifulSoup com o parser \texttt{lxml} para processar o HTML e produzir um texto limpo e compacto.

As tags removidas incluem: \texttt{script}, \texttt{style}, \texttt{noscript}, \texttt{iframe}, \texttt{svg}, \texttt{meta}, \texttt{link}, \texttt{head}, \texttt{nav}, \texttt{footer} e \texttt{aside}. Os atributos preservados são apenas aqueles que carregam informação semântica relevante: \texttt{alt}, \texttt{title}, \texttt{aria-label} e \texttt{href}. O texto é extraído com separadores de nova linha e linhas em branco consecutivas são reduzidas a uma única, reduzindo o tamanho do documento.

Para respeitar os limites de janela de contexto dos modelos e controlar o custo de processamento, o conteúdo é truncado a um máximo de 80.000 caracteres, configurável via variável de ambiente.

\subsection{Estágio 3: Extração (\textit{Extract})}

O estágio de extração é o núcleo diferencial do sistema. O texto limpo obtido no estágio anterior é enviado à API da OpenAI juntamente com dois elementos:

\begin{itemize}
    \item \textbf{Prompt de sistema}: instrui o modelo a atuar como um motor de extração de dados estruturados, retornando exclusivamente um objeto JSON válido, sem adicionar informações não presentes na página e sem seguir eventuais instruções encontradas no conteúdo da página (\textit{prompt injection});
    \item \textbf{Prompt de usuário}: contém o esquema JSON desejado — com os campos a serem extraídos e seus tipos — e o contexto opcional fornecido pelo usuário.
\end{itemize}

A chamada à API utiliza o parâmetro \texttt{response\_format: \{type: "json\_object"\}}, que garante que a resposta do modelo seja sempre um JSON válido. A temperatura é configurada como zero para maximizar o determinismo das respostas. O modelo padrão é o \texttt{gpt-4o-mini}, configurável via variável de ambiente, permitindo a substituição por modelos mais potentes ou mais econômicos conforme a necessidade.

\section{Tecnologias Utilizadas}

% TODO: Criar tabela com as tecnologias divididas por camada (Frontend, Worker, Infraestrutura).
% Incluir: React, TanStack Start, Bun, Python, BeautifulSoup, Playwright, OpenAI SDK,
% RabbitMQ, PostgreSQL, Drizzle ORM, Docker, Better Auth.

\section{Abordagem Tradicional: Módulo \textit{Baseline}}

% TODO: Descrever a implementação do módulo de scraping tradicional que será usado para comparação.
% Cobrir:
% - Estrutura do módulo (BeautifulSoup + seletores manuais)
% - Como os seletores foram definidos para cada cenário de teste
% - Critério para definição dos seletores (inspeção manual do DOM)
% - Diferenças de implementação em relação ao módulo LLM

\section{Metodologia de Comparação}

\subsection{Cenários de Teste}

% TODO: Descrever os cenários de teste selecionados.
% Proposta de categorias:
% 1. Páginas de produto em e-commerce (preço, título, descrição, disponibilidade)
% 2. Artigos de portal de notícias (título, autor, data, corpo do texto)
% 3. Páginas com dados tabulares (tabelas de dados, rankings)
% Definir o número de URLs por categoria e critérios de seleção.

\subsection{Métricas de Avaliação}

A comparação entre as duas abordagens será realizada com base nas seguintes métricas:

\begin{itemize}
    \item \textbf{Acurácia}: proporção de campos extraídos corretamente em relação ao total esperado, calculada por comparação com um gabarito construído manualmente;
    \item \textbf{Cobertura}: proporção de campos extraídos (independentemente de estarem corretos), indicando a completude da extração;
    \item \textbf{Tempo de processamento}: tempo total decorrido desde o recebimento da URL até a disponibilização do resultado, medido em segundos;
    \item \textbf{Custo por extração}: custo financeiro estimado por URL processada. Para a abordagem LLM, calculado com base no número de tokens consumidos e na tabela de preços da OpenAI. Para a abordagem tradicional, considerado como zero em custo de API;
    \item \textbf{Taxa de falha}: proporção de URLs que resultaram em erro ou sem dados úteis;
    \item \textbf{Esforço de implementação}: medido pelo tempo necessário para configurar o extrator para um novo site, avaliado qualitativamente.
\end{itemize}

% ============================================================
\chapter{Resultados e Discussão}
\label{cap:resultados}
% ============================================================

% TODO: Preencher este capítulo após a execução dos experimentos.

\section{Resultados da Abordagem LLM}

% TODO: Apresentar tabelas e gráficos com:
% - Acurácia por cenário de teste
% - Cobertura por cenário
% - Tempo médio de processamento
% - Custo médio por extração
% - Taxa de falha

\section{Resultados da Abordagem Tradicional}

% TODO: Apresentar tabelas e gráficos com as mesmas métricas para o baseline tradicional.

\section{Análise Comparativa}

% TODO: Comparar diretamente os resultados das duas abordagens:
% - Tabela consolidada com todas as métricas lado a lado
% - Gráficos comparativos
% - Identificar em quais cenários cada abordagem se saiu melhor

\section{Análise de Custo-Benefício}

% TODO: Discutir o trade-off entre custo financeiro (LLM) e custo de manutenção (tradicional).
% Propor uma análise de viabilidade econômica para diferentes volumes de scraping.

\section{Discussão}

% TODO: Interpretar os resultados à luz da revisão de literatura.
% Discutir limitações do estudo: número de sites testados, período de coleta, etc.
% Levantar hipóteses sobre os resultados inesperados.

% ============================================================
\chapter{Conclusão}
\label{cap:conclusao}
% ============================================================

% TODO: Escrever após a conclusão dos experimentos e análise dos resultados.
% Estrutura sugerida:

\section{Contribuições do Trabalho}

% TODO: Sintetizar as principais contribuições:
% - Sistema Crawler AI implementado e documentado
% - Comparação empírica formal entre abordagens de scraping
% - Evidências sobre quando usar LLM vs tradicional
% - Código-fonte disponível como referência

\section{Limitações}

% TODO: Discutir limitações reconhecidas:
% - Número limitado de sites testados
% - Dependência de API externa (OpenAI) sujeita a mudanças de preço
% - Comparação realizada em um momento específico (sites mudam)
% - Não avaliação de outros LLMs (Claude, Gemini, etc.)

\section{Trabalhos Futuros}

% TODO: Propor extensões ao trabalho:
% - Suporte a múltiplos provedores de LLM (Claude, Gemini, modelos locais como Llama)
% - Avaliação de técnicas de fine-tuning para a tarefa de extração
% - Desenvolvimento de um benchmark público para comparação de scrapers
% - Sistema de detecção automática da melhor abordagem por tipo de página

\bibliographystyle{abnt}
\bibliography{bibliografia}

\end{document}
